% conclusion.tex
\section{Conclusion}
\label{sec:conclusion}

We presented KnowFeat, a knowledge-guided feature engineering framework that combines hierarchical domain knowledge injection, multi-stage verification (L1--L3), and intelligent search strategies within an LLM agent. On seven public datasets, KnowFeat achieves the best AUC on 4 out of 7 and the best average rank (1.4) among all methods, while generating 3--10$\times$ fewer features than FeatLLM. On a real-world AML dataset with 200K transactions, 18 knowledge-guided features improve Recall by +2.1pp and F1 by +1.5pp over a strong 93-feature baseline---gains that are robust across different class balancing settings. The ablation study confirms that domain knowledge is essential, and every generated feature is traceable to specific knowledge assets.

\textbf{Limitations and Future Work.}
First, our current knowledge hierarchy is manually organized; automating knowledge extraction from heterogeneous sources (regulations, case law, expert interviews) is a promising direction.
Second, KnowFeat currently uses a single LLM; a multi-agent architecture with specialized agents for knowledge retrieval, feature proposal, and verification could improve efficiency.
Third, extending KnowFeat to additional high-stakes domains such as cybersecurity and insurance fraud detection would validate the generality of the approach.
Finally, investigating adaptive L3 thresholds to balance feature quantity and quality across different dataset characteristics is an important direction.
